%!TEX root = ../../main.tex
% ============================================================
% estrutura/pre-textuais/capa.tex
%
% DADOS DO TRABALHO — edite esta seção antes de tudo
% ============================================================

% ── Título e subtítulo ───────────────────────────────────────
\titulo{Infinitos de Cantor}
\subtitulo{Cardinalidade, Dimensão e a Estrutura dos Números Ordinais}

% ── Título em inglês (usado no Abstract) ────────────────────
\titleabstract{Cantor's Infinities: Cardinality, Dimension, and the Structure of Ordinal Numbers}

% ── Autor ────────────────────────────────────────────────────
\autor{Eduardo Furlan}
\autorcitacao{FURLAN, Eduardo}  % Sobrenome em maiúsculo

% ── Local e ano ──────────────────────────────────────────────
\local{Pato Branco}
\data{2026}

% ── Natureza do trabalho ─────────────────────────────────────
% Opções: Trabalho de Conclusão de Curso | Dissertação | Tese
%         Projeto de Qualificação
\projeto{Trabalho de Conclusão de Curso}

% ── Título acadêmico ─────────────────────────────────────────
% Opções: Bacharel | Tecnólogo | Licenciado | Mestre | Doutor
\tituloAcademico{Licenciado}

% ── Instituição ──────────────────────────────────────────────
% Se for TCC, coloque o nome do curso em \programa{}
% \logoinstituicao{<escala>}{<caminho/nome do arquivo>}
\instituicao{\textbf{Universidade Tecnológica Federal do Paraná}}
\departamento{\textbf{DAMAT}}
\programa{\textbf{Curso de Licenciatura em Matemática}}
\logoinstituicao{0.2}{../../dados/figuras/logo-instituicao.png}

% ── Área de concentração / linha de pesquisa ─────────────────
% Deixe vazios para TCC de graduação
\areaconcentracao{}
\linhapesquisa{}

% ── Orientador(a) ────────────────────────────────────────────
\orientador{Prof. Dr. Ivan Italo Gonzales Gargate}
% \orientador[Orientadora:]{Profa. Dra. Nome da Orientadora}
\instOrientador{Universidade Tecnológica Federal do Paraná}

% ── Coorientador(a) — descomente se houver ──────────────────
% \coorientador{Prof. Dr. Nome do Coorientador}
% \instCoorientador{Instituição do Coorientador}
